\section{Discussion}

\subsection{The Secret Use of LLMs as a Prevalent Challenge for AI Transparency}

The findings of our study suggest that a large portion of the populations generally tend to hide the use of LLMs for various tasks (\autoref{sec:desc-stats-secret-use-estimates}).
This is a concerning issue, as we observe a significant overlap between the tasks in which users exhibited the secretive behaviors, and the tasks in critical domains such as academic writing, school work, creative writing, and professional work tasks.
As we have reviewed in \autoref{sec:community-level-requirements-of-transparency}, these domains have regulations or guidelines that explicitly require a disclosure of how LLMs are employed in achieving the tasks.
While having these regulations is the first step towards establishing clear rules and expectations for the transparency requirements of AI users, our study results suggest that non-compliance with these regulations may be prevalent.


When the use of LLM in critical tasks is not transparently disclosed, potential issues in the content generated by LLM, such as trade secret, bias or misinformation, can spread widely without sufficient precautionary measures being taken.
For example, using AI to research a court case has led to unwittingly including fictitious legal citations in a court brief~\cite{merken2023new}. 
In fact, our participants have shown some level of awareness of the potential negative consequences of their behavior. 
For example, EP80 (P80 in Study 2 - Experiment) thought~\blockquote{Using AI for detailed market research reports makes it feel both unreliable in a way.}, yet she still chose to keep her usage secret in both passive and active way. 
The seemingly contradictory behavior imply the existing challenges in ensuring that individual users do their part to provide the necessary transparency to prevent, mitigate, and diagnose the risks caused by AI-generated content.


Our findings regarding the mechanisms of users' intentions on the secret use of LLMs suggest that promoting transparency in LLM use cannot be a one-size-fits-all approach. 
Interventions must balance the need for AI transparency with individual privacy concerns, considering the context and varying levels of secretive behavior. 
Our findings emphasized the important role of users' perceived external judgment (perceived social norms) in shaping these behaviors and calls for tailored strategies to encourage responsible, transparent AI use by end-users.
\revision{Although we did not observe significant mediation effects from internal judgment as a whole in the quantitative analysis, qualitative findings from our first study did reveal some impact of users' internal judgment on their intentions to secret use, such as how using LLMs in a task affects their perceived self-efficacy and how it aligns with their own moral standards.
Our controlled experiment focused only on the mediation effect of two types of user perceptions towards LLM usage on their intention to secret use, yet there might be more complexities within the underlying mechanism.
Future research can look deeper into the mechanism of users' decision-making process in secret use based on our initial findings.
For example, including users' actual willingness to use LLMs in specific scenarios as a control variable can potentially provide a more comprehensive understanding of the decision-making process.}



\subsection{Addressing Tensions between Transparency Needs and Individual Privacy Concerns}


In our study, multiple tasks engendered high intentions for hiding the use of LLMs, while the potential harm levels associated with the lack of transparency varies.
Hence, we argue that it is beneficial to employ a multi-layered approach to encourage transparent disclosure of AI use among end-users.
In situations with less severe consequences, persuasive interventions like community policies and behaviorial nudges can guide or encourage transparent AI usage. 
For example, video platforms like YouTube and TikTok developed an AI labeling system to encourage creators to disclose their AI tool usage in video creation to audience~\cite{YouTube2024ResponsibleAI, TikTok2024AIGenerated}.
Our findings about the significant role of perceived external judgements suggests that social norm nudges~\cite{bicchieri2023norm} may be particularly relevant in addressing the users' concerns that lead to concealment intentions.
In addition, positive incentives such as a reward policy, could be more effective instead of only negative inhibition such as warning or punishment policy. 
This presents a promising direction for future research.
On the other hand, in critical areas where irresponsible usage could have severe consequences, or by adversarial users, soft nudges may not be enough. 
Therefore, in such high-stakes scenarios, stricter measures such as monitoring AI use, detecting AI-generated artifact, and implementing comprehensive policies and laws might be necessary.



\subsection{Emotional Stress Associated with the Secret Use of LLMs}
We want to note emotional stress as a potential negative impact of the secret use behavior on the user themselves.
Apart from contexts with more impact on the public domains, contexts like ``social connection'', where LLMs are used for personal messaging, caused the highest concealment intentions (92\% for passive non-disclosure; 70\% for active concealment). 
Qualitative responses from our survey and experiment suggest that users often experience emotional stress in these situations. 
For example, SP40 used ChatGPT for writing love letters to her spouse while keeping it secret. 
She said, ~\blockquote{...I felt guilty. My partner thought that I had said all of those kind words, but I did not}.
EP190 used an LLM tool to write a thank-you letter for a gift while she revealed bad feelings for herself,~\blockquote{I really did not feel involved enough to do so. It takes away from what it meant to write something to them if they were to know.}
In addition to the qualitative evidence, theoretical frameworks such as the self discrepancy theory also suggest that a conflict between users' self-beliefs about using LLMs for social messaging and their actual behavior may lead to self-discrepancies and subsequent emotional stress~\cite{carver1999self, scott1993self}.
Future research needs to be conducted to examine the issue of emotional stress in terms of how it may lead to changes in behaviors and long-term well-being.



\subsection{Towards the Establishment of Community Norms of AI Use and Disclosure}

Perceived external judgement on the use of LLMs, which was confirmed to be an important factor in shaping users' intentions of concealment of the usage, is a kind of perceived social norms towards this technology use~\cite{choi2013applying, higgs2015social}.
We are currently at the stage where LLMs emerge and open up new fields of human-AI interaction in various aspects of life.
Nevertheless, there remain uncertainties about how well users' perceptions align with actual social norms~\cite{diefenbach2019disrespectful,hechter2001have}.
If the actual reactions to LLM usage are positive or neutral, contrary to users' negative perceptions, it presents an opportunity to promote transparency by aligning users' perceived external judgments with the actual ones.
Therefore, we advocate going beyond just developing standards in a top-down manner, and emphasize the importance of engaging with a bottom-up approach to investigate people's perceptions of AI use in specific communities.
As an example, \citet{kapania2024m} conducted surveys and interview studies with HCI researchers to examine their perceptions of the ethical issues of LLM use in HCI research practices including ideation and project scoping, study design and execution, and analysis and paper writing.
These types of studies will play an instrumental role in understanding the community norms surrounding AI use.

\revision{Moreover, future research may need to look further into developing support for creating standardized, accurate, and meaningful disclosure of LLM use by systematically reviewing factors such as cultural differences, use domains, or varying levels of LLM operational involvement.
For example, the use of LLMs in personal vs. professional domains (e.g., maintaining personal relationships v.s. making professional medical decisions) or at different levels of involvement (e.g., refining language based on user drafts vs. generating the entire content) could potentially lead to different types or levels of risks.
The ideas of designing standardized, concise, and machine-readable privacy nutrition labels~\cite{kelley2009nutrition, kelley2010standardizing, emami2020ask}, as well as tools that make it easy to create these labels~\cite{li2022understanding, li2024matcha}, serve as a good example of addressing the transparency issues in data privacy. These concepts can be inspirational in addressing the transparency issues related to the secret use of LLMs.
Overall, we consider the design of support to facilitate community-level deliberation of norms and best practices of usage disclosure by AI users as a crucial and nascent area for future research.}